\section{Partial turns and the terminal wall}\label{sec:partial-turn}

Put $L=\pi/2$. In this section the whole terminal wall is retained:
\[
 \begin{aligned}
 N_\alpha(U)(x)&=\max\{n_\alpha(U)(x),a_\alpha(x)\},\\
 n_\alpha(U)(x)&=\max\{0,\sup_{0<t<\alpha}\min(a_t(x),b_t(x))\}.
 \end{aligned}
\]
The two functions agree at $\alpha=L$, because $a_L\le0$.
Horizontal reflection exchanges the two wall families. At a proper
angle it does not preserve the objective with its first terminal wall.
All changes of orientation below concern a symmetric finite-angle
comparison or the support equations, and keep this distinction.

\subsection{The remaining geometry of a proper maximum}

Fix a joint maximizing pair $(U,\alpha)$, with $\alpha<L$.
The canonical reductions of \cref{lem:selection-canonical} keep its
own angle and score. Center its projection and write
\[
 \begin{gathered}
 I_U=[-2C,2C],\quad J_U=[-C,C],\\
 f_U(-C)=1-h_L,\quad f_U(C)=1-h_R,\\
 h_L,h_R\ge0,\quad\min(h_L,h_R)=0.
 \end{gathered}
\]
The saturated description in \cref{lem:selection-saturation} is used
throughout. In particular an unused normal interval has zero charged
curvature, apart from a possible uncharged middle facet.

\begin{lemma}[Angle and width cuts]\label{lem:partial-cuts}
For a canonical joint maximum at a proper angle,
\[
 \alpha>\arctan(8/3),\qquad
 \frac{1001}{2000}<C<\frac{37}{50},\qquad
 h_L+h_R<\frac{17}{50}.
\]
\end{lemma}
\begin{proof}
The reference cap is a competitor, so the maximizing score is at
least $\Mc/2$. The initial-angle estimate in
\cref{app:partial-initial} is strictly below $\Mc/2$ for
$\pi/4\le\alpha\le\arctan(8/3)$. All three angles
$\pi/8,\pi/4,3\pi/8$ have therefore been visited.
The finite-angle roof and triangle estimates in
\cref{lem:preliminary-cuts} give $1001/2000<C<4/5$ and
$h_L+h_R<17/50$.
They apply to the partial barrier because it dominates the same
three tents. If the high middle endpoint is on the left, reflect
the cap and its actual barrier together in this comparison.
The endpoint-pressure equations are merely relabeled; the outgoing
wall is not changed to the wall of another variational problem.
The revised floor-loss calculation in \cref{app:partial-width}
then excludes $C\ge37/50$.
\end{proof}

\begin{lemma}[Completion of a descending middle]\label{lem:partial-completion}
Suppose $h_R=h>0$ and put
$\theta_0=L-\arctan(h/(2C))$. If
$\alpha\ge\theta_0$ or $h\ge1-\sin\alpha$, then
$N_\alpha(U)\ge n_L(U)$ on $J_U$ and
$P_\alpha(U)\le P_L(U)$.
Otherwise $0<h<1-\sin\alpha$ and $\alpha<\theta_0$.
\end{lemma}
\begin{proof}
The negative affine middle has its top endpoint at $(-C,1)$.
Saturation gives, on the missing first-normal arc,
\[
 h_U(t)=
 \begin{cases}
 C\cos t+(1-h)\sin t,&\alpha<t<\theta_0,\\
 -C\cos t+\sin t,&\theta_0<t<L.
 \end{cases}
\]
The first case is omitted when $\alpha\ge\theta_0$.
No charged extension of the middle facet into an unused normal arc
is possible. The second branch has first wall
$(-C-x)\cot t+1-\csc t\le0$ on $J_U$.
On the first branch put $z=C-x\in[0,2C]$; its wall is
\[
 A_z(t)=1-h+\frac{z\cos t-1}{\sin t},\qquad
 A_z'(t)=\frac{\cos t-z}{\sin^2t}.
\]
For $z\ge\cos\alpha$ it decreases after $\alpha$, and is at
most the retained terminal wall. For $0\le z<\cos\alpha$ its
maximum on $[\alpha,L]$ is
$1-h-\sqrt{1-z^2}\le1-h-\sin\alpha$.
Under the asserted height condition this is nonpositive. Every
missing two-wall minimum is consequently bounded by
$\max(0,a_\alpha)$, proving the pointwise comparison and the
score inequality.
\end{proof}

Discard the completed descending branch for the moment. In every
other case set
\begin{equation}\label{eq:partial-tail-data}
 c=\cos\alpha,\quad s=\sin\alpha,\quad H=1-h_R,
 \quad d=\frac{1-Hs}{c},\quad w=c-d.
\end{equation}
Then $s<H\le1$, $0<d<c$ and $0<w<c/2$.
For a horizontal or rising middle $H=1$; for a descending one
these inequalities follow from the residual alternative in
\cref{lem:partial-completion}.
Let $T_L,T_R$ be the height-one top overhangs outside $J_U$.
They vanish on a strictly lower middle side. Immediately after
$\alpha$ the first support is the point $(b,H)$, where
$b=C+T_R$; in the descending case this remains true up to
$\theta_0$. The companion support after $\alpha$ is the left
endpoint $(-2C,e_L)$, and has no terminal atom.
If $\tau$ is the horizontal length of the first terminal facet,
its endpoints and shifted endpoints are
\begin{equation}\label{eq:partial-terminal-facet}
 (b,H),\ (b+\tau,H-\tau c/s),\qquad
 B_+=(b-c,H-s),\quad B_-=B_++(\tau,-\tau c/s).
\end{equation}
In particular
\begin{equation}\label{eq:partial-terminal-wall}
 a_\alpha(-C+z)=\frac cs(2C+T_R-d-z),\qquad
 r_0=b-d
\end{equation}
is the terminal wall and its zero.

\begin{lemma}[Top projection]\label{lem:partial-top-projection}
For the residual geometry above, $T_L+2T_R\le d$ and $T_R<d$.
\end{lemma}
\begin{proof}
By \cref{cor:variation-projection},
\[
 2C-T_L-T_R\ge|\{x\in J_U:N_\alpha(U)(x)>0\}|
       \ge|J_U\cap(-\infty,r_0)|.
\]
If $T_R\ge d$, the last interval has length $2C$, forcing
$T_L=T_R=0$, contrary to $T_R\ge d>0$.
Thus $T_R<d$. Also $r_0>C-c>0$, and the last interval has
length $2C+T_R-d$. Rearrangement proves the result.
\end{proof}

\subsection{One-sided angle variations and the terminal facet}

\begin{lemma}[The angle derivative with ties]\label{lem:partial-angle-derivative}
Let $E=\{x\in J_U:a_\alpha(x)>n_\alpha(U)(x)\}$ and
$Z=\{x\in J_U:a_\alpha(x)=n_\alpha(U)(x)\}$.
At a proper angle the right derivative, with the cap fixed, is
\[
 \partial_+P_\alpha(U)
 =-\frac1{s^2}\left[
       \int_E(x-(B_+)_x)\,\dd x
       +\int_Z(x-(B_+)_x)_+\,\dd x\right].
\]
At a joint maximum the bracket is nonnegative.
\end{lemma}
\begin{proof}
The right support derivative chooses the upper/left endpoint of
the supporting face. Uniformly on $J_U$,
$a_{\alpha+\varepsilon}(x)=a_\alpha(x)
+\varepsilon(x-(B_+)_x)/s^2+o(\varepsilon)$.
If $a_\alpha>n_\alpha$, the terminal companion wall is strictly
lower, since otherwise its limit from earlier angles would already
give $n_\alpha\ge a_\alpha$. Newly visited corners therefore
cannot alter the first-order terminal derivative.
If $a_\alpha<n_\alpha$, the old history supplies a strict gap,
giving derivative zero. On a tie,
\[
 \max(n_\alpha,a_{\alpha+\varepsilon})\le
 N_{\alpha+\varepsilon}\le
 \max(n_\alpha,\sup_{0\le t\le\varepsilon}a_{\alpha+t});
\]
both bounds have derivative $((x-(B_+)_x)/s^2)_+$.
The supports have locally bounded one-sided derivatives, and the
angle denominators are bounded away from zero. The quotients are
uniformly bounded on $J_U$, so dominated convergence gives the
integrated formula. Its sign follows from maximality.
\end{proof}

\begin{lemma}[A short facet for each prescribed maximum]\label{lem:partial-facet}
Every residual canonical proper-angle joint maximum satisfies
$\tau\le w$. No largest-angle choice is required.
\end{lemma}
\begin{proof}
The terminal corner $z_\alpha$ satisfies
\[
 (z_\alpha)_x-(B_+)_x
 =s[1-(b+2C)s+(H-e_L)c]<0.
\]
Indeed $b\ge C>1/2$, $H\le1$, $e_L\ge0$, and
$s>12/13,c<5/13$, so $3Cs-c>1$.
Suppose $\tau>w$. Then $(B_-)_x>r_0$.
For $(z_\alpha)_x<x<(B_-)_x$, the companion terminal wall
exceeds the first one, while
$\partial_-a_\alpha(x)=(x-(B_-)_x)/s^2<0$.
The companion has no terminal atom. Sufficiently close earlier
angles have both walls above $a_\alpha(x)$, and hence
$n_\alpha(U)(x)>a_\alpha(x)$ on this interval.
For $x\ge(B_-)_x$ the first wall is negative.
Thus strict terminal exposure and relevant ties lie to the left
of $(z_\alpha)_x<(B_+)_x$.
\Cref{lem:partial-angle-derivative} forces $|E|=0$.
Continuity then gives $N_\alpha=n_\alpha$ on $J_U$.

Choose
\[
 \max(-C,(z_\alpha)_x)<\zeta<(B_+)_x<r_0
       <r_*<\min(C,(B_-)_x).
\]
Immediately after $\alpha$ the same cap has first support $(b,H)$,
so $\partial_ta_t(x)=(x-b+\cos t)/\sin^2t$.
For $x\le\zeta$ its first walls decrease for a small angle
increment. On $[\zeta,r_*]$ the old history has a uniformly
positive gap above $a_\alpha$; uniform convergence covers the
new walls there. The new zero
$b-(1-H\sin t)/\cos t$ remains below $r_*$, so the new first
walls are nonpositive to its right. Choose the increment below
$\theta_0-\alpha$ when the middle descends.
For some $\beta>\alpha$, all new first walls are consequently
below $n_\alpha$ on $J_U$. Their two-wall minima are below the
same bound. The old history remains present, and
$N_\beta=N_\alpha$, so $(U,\beta)$ is another joint maximum.

The cap and its window have remained fixed. Its saturation persists:
$U\subseteq U^{\mathrm{sat}}_\beta\cap\{y\le\ell_U\}
\subseteq U^{\mathrm{sat}}_\alpha\cap\{y\le\ell_U\}=U$.
The original terminal facet has positive charged length and normal
$\alpha\in(0,\beta)$, away from the middle normal.
\Cref{thm:variation-curvature} for this prescribed new maximizing
pair excludes such an interior charged atom, a contradiction.
\end{proof}

\begin{lemma}[Strict terminal exposure]\label{lem:partial-exposure}
For a canonical proper-angle maximum, $|E|\le\tau$.
\end{lemma}
\begin{proof}
Move its terminal supporting line inward by $\varepsilon$.
The endpoint baseline points remain present, so the projection and
the window do not change. The exterior area lost to first order is
$\varepsilon\tau/s$: the charged facet has length $\tau/s$.
At every strict terminal contact the barrier drops to first order
by $\varepsilon/s$. Every other used support can only decrease,
which can only improve the score. Restrict first to a compact subset
of $E$ with a positive contact gap and then exhaust $E$.
The inward derivative of the score is at least
$(|E|-\tau)/s$. Maximality makes it nonpositive, proving the
inequality. At a zero-length facet the exterior loss is
$o(\varepsilon)$ and gives the same conclusion.
\end{proof}

\subsection{A weighted companion comparison}

For a rising middle use the left charged-wing support $G$ alone.
The omitted right-wing points have $X\ge C,Y\le1$, and their
companion walls on $J_U$ are at most
$1-\sec t+(x-C)\tan t\le0$. The middle segment is controlled
by its endpoints. This replacement therefore preserves positive
niche contacts. For the other signs use the actual companion
support. In all cases
$G(0)=1-h_L$, $G'(0+)=C+T_L$ and $v=G''+G\ge0$.
Put
\begin{equation}\label{eq:partial-weighted-error}
 \begin{gathered}
 A_0=C-T_L,\quad\theta=\arcsin A_0,\quad\lambda=\tan\theta,\\
 \mathcal E_\theta=\int_0^\theta(v(t)-1)_+
                 (\lambda\cos t-\sin t)\,\dd t.
 \end{gathered}
\end{equation}

\begin{lemma}[Localized companion moment]\label{lem:partial-moment}
For $0\le z\le w$,
\[
 n_\alpha(U)(-C+z)\le
 \sqrt{1-(C-T_L)^2}-\sqrt{1-(C+z)^2}+\mathcal E_\theta.
\]
\end{lemma}
\begin{proof}
Write the shifted companion support point as
$D=(G-1)u^\perp-G'u$. Then
\[
 D_x'=(1-v)\cos t,\quad D_y'=(1-v)\sin t,\quad
 b_t'(x)=\frac{x-D_x(t)}{\cos^2t}.
\]
Its outer abscissa is $D_x-\sin t\ge-2C$.
Also $0<A_0<C+w<C+c/2<s$; the final inequality follows from
$C<37/50,c<3/\sqrt{73}$ by squaring
$37/50+3/(2\sqrt{73})<8/\sqrt{73}$.
A positive maximum of the companion-wall family on $[0,\alpha]$
is interior: its value at zero is $-h_L$, and its derivative is
negative once $\sin t>C+z$. At a maximizing angle
$D_x=-C+z$, $b_t=D_y$ and $t\le\arcsin(C+z)$.
Absolute continuity of the regular support gives
\[
 D_y(t)+h_L-\lambda(T_L+z)
   =\int_0^t(1-v(r))(\sin r-\lambda\cos r)\,\dd r.
\]
Before $\theta$ its integrand is bounded above by the weighted
excess. After $\theta$ it is at most
$\sin r-\lambda\cos r$ because $v\ge0$.
For $t\ge\theta$, integration bounds $D_y+h_L$ by
\[
 \sqrt{1-A_0^2}-\cos t
       +\lambda(C+z-\sin t)+\mathcal E_\theta.
\]
The part preceding the excess increases with $t\ge\theta$,
and its maximum at $t=\arcsin(C+z)$ is the asserted circle
difference. For $t<\theta$, use instead
$D_y+h_L\le\lambda(T_L+z)+\mathcal E_\theta$ and
\[
 \lambda(T_L+z)\le
 \int_{A_0}^{C+z}\frac{x}{\sqrt{1-x^2}}\,\dd x.
\]
Finally the niche is bounded by the positive part of the largest
companion wall. The asserted right side is nonnegative, and
discarding $-h_L$ proves the result even when there is no positive
wall maximum.
\end{proof}

\begin{lemma}[The terminal contradiction]\label{lem:partial-tail-contradiction}
No canonical joint maximum has proper angle and residual tail geometry.
\end{lemma}
\begin{proof}
The exact circle comparison in \cref{app:partial-circle} gives
\[
 a_\alpha(-C+z)-
 \left[\sqrt{1-(C-T_L)^2}-\sqrt{1-(C+z)^2}\right]
       >\frac{39}{4400}c\qquad(0\le z\le w).
\]
The reflected support and terminal-impulse estimates in
\cref{app:partial-reflection} give $\mathcal E_\theta=0$ for
$0<\cot\alpha\le1/8$, and
\[
 \mathcal E_\theta<\frac{6517}{6400000}
       <\frac{117}{110000}<\frac{39}{4400}c
\]
for $1/8\le\cot\alpha<3/8$.
Together with \cref{lem:partial-moment} these inequalities make
the entire interval $[-C,-C+w]$ strictly terminal-exposed.
Its right endpoint is interior to $J_U$, so continuity extends
exposure a positive distance farther right. Thus
$|E|>w\ge\tau$, contradicting \cref{lem:partial-exposure}.
\end{proof}

\subsection{The sharp theorem for every cap and its own angle}

\begin{theorem}[Sharp cap theorem]\label{thm:cap-bound}\label{thm:cap-equality}
Every downward cap $U$ and every $\alpha\in[\pi/4,L]$ satisfy
\[
 P_\alpha(U)\le\Mc/2.
\]
Equality holds if and only if $\alpha=L$ and
$U=\KR+(t,0)$ for some $t\in\R$.
\end{theorem}
\begin{proof}
Choose a joint maximum with largest terminal angle, whose existence
is given by compactness. Canonicalization retains that angle.
If it is proper, \cref{lem:partial-cuts} applies. In a completed
descending branch, \cref{lem:partial-completion} gives
$P_\alpha(U)\le P_L(U)\le\Mc/2$ by
\cref{thm:full-cap-bound}. Since the reference is a competitor,
both inequalities are equalities and the same cap at $L$ is a
joint maximum, contrary to the largest-angle choice. Every other
proper branch is excluded by \cref{lem:partial-tail-contradiction}.
Thus the joint value is the full-angle value $\Mc/2$, proving
the universal bound.

Now let the \emph{prescribed} original pair $(U,\alpha)$ have score
$\Mc/2$. The bound makes this pair a joint maximum.
Its canonical replacement keeps its own angle and score.
If its angle is proper, the completed descending branch would make
that replacement a full maximum. Full-angle affine rigidity forces
its middle slope to vanish, contradicting descent.
Every residual branch is excluded by the same terminal contradiction;
\cref{lem:partial-facet} applies to this individual pair, without
any largest-angle selection. Hence its angle is $L$.
The original full-angle equality classification in
\cref{thm:full-cap-bound} then identifies $U$ itself as a
horizontal reference translate. The converse follows from the
reference full score and horizontal translation invariance.
\end{proof}
